\documentclass[a4paper]{article}
\usepackage{ctex}
\usepackage{geometry}
\usepackage{amsmath}
\usepackage{ulem}
\usepackage{graphicx}
\usepackage{siunitx}
\usepackage{enumitem}
\usepackage{multirow}
\usepackage{adjustbox}
\usepackage[hidelinks]{hyperref}

\geometry{left=2.5cm, right=2.5cm, top=2.5cm, bottom=2.5cm}
\zihao{-4}

\begin{document}

\noindent\textbf{班级：24级整合科学班} \hspace{1cm} \textbf{学号：2400017802} \hspace{1cm} \textbf{姓名：冯翊展} \hspace{1cm} \textbf{评分：}\\
\textbf{同组学生：胡家豪 \ 陈天铭 \ 侯树宽} \\
\textbf{实验名称：观察大肠杆菌鞭毛的转动}   \hfill \textbf{实验日期：2025年3月21/28日} 

\hrule
\vspace{3mm}

\begin{enumerate}[label=\chinese*.,left=0pt]
    \item 实验目的
    \\
    通过两种⽅式（Tethered cell和Bead assay）观察⼤肠杆菌鞭⽑的转动。
    
    \item 实验内容
    \begin{enumerate}[label=\arabic*.,left = 0pt]
        \item Tether cell：将⼤肠杆菌鞭⽑固定在盖玻⽚上，在显微镜下观察其转动情况。
        \item Bead assay：将大肠杆菌自身固定在盖玻片上，加入荧光微球使其粘附在鞭毛上，在荧光显微镜下观察小球的转动情况。
    \end{enumerate}
    
    \item 基本原理
    \begin{enumerate}[label=\arabic*.,left = 0pt]
        \item ⼤肠杆菌细胞膜上存在若⼲细长弯曲、具有运动功能的蛋白质附属丝状物，称为鞭毛。鞭毛的根部存在鞭毛马达，可将化学能转化为机械能，其转动可以控制鞭毛丝的摆动，从而控制细菌的运动。
        \item Tether cell：大肠杆菌在流道中沉降后鞭⽑会粘接在盖玻⽚上，翻转切⽚并在显微镜下观察，可以看到菌体在鞭⽑⻢达的驱动下做与鞭⽑转动⽅向相反的转动。
        \item Bead assay：⽤注射器削短⼤肠杆菌鞭⽑⻓度，加入Polylysine后菌体被固定在盖玻⽚上。再加⼊荧光珠，让其吸附在削短的鞭⽑上，便可以在荧光显微镜下观察荧光珠的转动，间接观察鞭⽑转动情况。
    \end{enumerate}
    
    \item 仪器与试剂
    \begin{enumerate}[label = \arabic*., left = 0pt]
        \item 仪器：\qty{100}{\micro \L}/\qty{1000}{\micro \L} 移液枪，离心管，离心机，注射器，载玻片，盖玻片，荧光显微镜，双⾯胶、吸⽔纸
        \item 试剂：LB培养基（酵母提取物、胰蛋白胨、氯化钠），TB培养基（胰蛋白胨、酵母提取物、甘油、磷酸一/二氢钾），流度缓冲液MB（\qty{10}{\milli M}磷酸钾、\qty{0.1}{\milli M}EDTA、pH=7），大肠杆菌鞭毛粘性突变体(\textit{E.coli.} SYC12），0.01\%聚赖氨酸（Polylysine），\qty{0.2}{\milli \m}与 \qty{1}{\milli \m}荧光微球溶液，指甲油
    \end{enumerate}

    \item 实验步骤
    \begin{enumerate}[label = \arabic*., left = 0pt]
        \item 细菌的准备
        \begin{enumerate}[label = 1.\arabic*., left =0pt]
            \item 实验前一天晚上，将\textit{E.coli.} SYC12以1:1000的比例转接于LB中，过夜活化。
            \item 次日9:30，将活化的菌1:100转接于TB中，在 \qty{30}{\degreeCelsius}振荡培养 \qty{3}{\hour} \qty{20}{\min}。
            \item 将培养好的细菌取出，用于实验。
        \end{enumerate}
        \item 实验一Tether cell
        \begin{enumerate}[label = 2.\arabic*., left =0pt]
            \item 搭建流道。
            \item 取 \qty{50}{\ul} - \qty{100}{\ul}的菌液，注入流道一端。
            \item 盖玻片向下静置\qty{10}{\minute}。
            \item 在流道一端滴加MB，另一端小心地用吸水纸吸，冲洗未黏附的细菌。
            \item 指甲油封片，上镜观察。
        \end{enumerate}
        \item 实验二Bead Assay
        \begin{enumerate}[label = 3.\arabic*., left =0pt]
            \item 取出 \qty{3000}{\g}菌液，离心 \qty{2}{\min}，移去TB，加入MB。
            \item 用 \qty{1}{\ml}注射器，反复吹打30次，削短鞭毛。
            \item 制作流道：用双面胶制备简易流道 （整个过程中盖玻片向下放置）。
            \item 向流道内注入Polylysine，孵育 \qty{5}{\min}。
            \item 用MB将Polylysine洗去，然后向流道中注入准备好的菌液，孵育 \qty{10}{\min}。
            \item MB洗去未粘上的菌，然后注入珠子（30倍稀释于MB中），孵育 \qty{10}{\min}。
            \item MB洗去未黏上的珠子（此步倾斜玻片轻柔洗去）。
            \item 上镜观察。

        \end{enumerate}
    \end{enumerate}

    
    \item 实验结果与分析
    \begin{enumerate}[label = \arabic*., left = 0pt]
        \item Tethered cell：
        \begin{enumerate}[label = 1.\arabic*., left =0pt]
            \item 部分菌体静⽌不动，可能是由于在培养、处理过程中细菌失去活性，进⼊休眠状态或已死亡，导致鞭⽑不再旋转。
            \item 作旋转的菌体中大部分都在作顺时针旋转，这表明观察到的大肠杆菌鞭毛在作逆时针旋转，即处在鞭毛束合并，推动细菌直线“游动”（run）的状态，这符合大肠杆菌在流道中趋向性的运动状态。
            \item 大肠杆菌的旋转中心有多种分布，可能出现在细菌的一端或者接近中心的位置，这表明大肠杆菌表面存在多根鞭毛可以旋转。
            \item 大肠杆菌的旋转速度不一，有的仅有 \qty{0.2}{rpm}，有的高达 \qty{5}{rpm}，这表明大肠杆菌鞭毛的转速受到环境阻力影响较大            
        \end{enumerate}
        \item Bead assay:
        \\
        相较前一个方法，本方法观测到的小球大部分均为静止状态，即使是运动的小球大部分也处于无规运动中，仅有少部分荧光小球在进行有规律的旋转，可能的原因是
        \begin{enumerate}
            \item 在盖玻片向下静置的过程中，大部分细菌整体粘附在了盖玻片上，导致无法发生旋转。
            \item Polylysine本身粘附了大量荧光小球，导致实际观察中大部分荧光小球都没有移动。
            \item 使用MB冲洗菌体和荧光小球的力度和次数难以控制，次数太多容易冲走粘附的细菌，次数太少则会使游离的荧光小球过多。
            
        \end{enumerate}

    \end{enumerate}



\end{enumerate}
\end{document}